\chapter{Fourier 变换}\label{ch:5}\index{Fourier–Deligne 变换|hyperpage}
% 原PDF第32页下半部
\section{回顾：从 Laumon 层到初步的自守层}\label{sec:5.1}
继续取 $G=\GL_n$，并令 $E$ 为 $X$ 上秩 $n$ 的局部系统。\chref{4}构造了 $\cL_E\in\operatorname{Perv}(\Tor)$。这里的目标是用它在 $\Bun_n$ 上构造自守层 $\Aut_E$。

令 $\Bun'_n$ 为对 $(M,s)$ 的叠，其中 $M$ 是秩 $n$ 向量丛，$s:\Omega_X^{n-1}\hookrightarrow M$ 为单射。令 $Q_n$ 为正的整体 Hecke–Whittaker 叠
\[Q_n=\left\{\begin{array}{l}0=M_0\subset M_1\subset\cdots\subset M_n,\quad M_i/M_{i-1}\simeq\Omega_X^{n-i},\\M_n\hookrightarrow M\in C_n\text{，余核为挠层}\end{array}\right\}.\]
等价地，这是与 $\rho$ 扭曲 Whittaker 叠作纤维化的正 Hecke 对应的紧化整体版本；它不是局部仿射 Grassmann 簇与 $\Bun_N$ 的字面乘积。它有三个显然的映射：$\pi:Q_n\to\Bun'_n$ 记住 $\Omega_X^{n-1}=M_1\hookrightarrow M$；$p_1:Q_n\to\Tor$ 记录挠商 $M/M_n$；$p_2:Q_n\to\Bun_N\to\Ga$ 是 Whittaker 特征标，即旗中相邻扩张类的和。
\begin{transnote}{含糊点拨}
「$\rho$ 扭曲」指把 $\Bun_N$ 上的 Whittaker 层按半和正根 $\rho$ 作 Tate 扭转：沿 Hecke 修改搬运 $N$-丛时，次数按 $\langle\rho,\cdot\rangle$ 平移，不扭则 $p_2^*\cL_\psi$ 与 $\pi$、$p_1$ 的等变性不能同时成立。称其为「紧化整体版本」，是指旗 $0=M_0\subset\cdots\subset M_n$、挠商 $M/M_n$ 与 Whittaker 特征标捆在同一叠里：末项 $M_n\hookrightarrow M$ 由同一条链给出，不能拆成 $\Gr\times\Bun_N$ 的直积；被丢掉的正是有挠缺陷（$M/M_n\ne 0$）的修改，而那恰是 $j_{k,!}$ 要延拓过去的部分。[依据：原书 p.33 本段；并见 p.22 §3.1 的 $\rho$-移位。]
\end{transnote}

令 $\cL_\psi$ 为 $\Ga$ 上的 Artin–Schreier 层。在 $\length(M/M_n)=d$ 的分支上，初步对象为
\[\Aut_E^{\prime,d}:=\pi_!(p_1^*\cL_E^d\otimes p_2^*\cL_\psi)[\dim_{\mathrm{rel}}](\dim_{\mathrm{rel}}/2).\]
方括号及圆括号表示 Fourier 归一化所用的相对移位与 Tate 扭转。感叹号表示紧支撑推前；映射 $\pi$ 一般不是真态射。

所希望的性质如下。
\begin{enumerate}
\item 若 $E$ 不可约，则对 $d\gg0$，$\Aut_E^{\prime,d}$ 是不可约反常层。
\item 对 $d\gg0$，$\Aut_E^{\prime,d}$ 沿 $\Bun'_n\to\Bun_n^d$ 下降，得到 $\Aut_E^d\in\operatorname{Perv}(\Bun_n^d)$。
\srcpage{33}
\item $\Aut_E$ 满足截断的最高次 Hecke 关系。采用\chref{2}的 Hecke 对应约定，
\[t_x^*\Aut_E^d\simeq\det(E_x)\otimes\Aut_E^{d+n},\]
其中 $t_x:\Bun_n^{d+n}\xrightarrow{\sim}\Bun_n^d$，$M\mapsto M(-x)$。等价地，$(t_x)_!\Aut_E^{d+n}\simeq\det(E_x)^{-1}\otimes\Aut_E^d$。
\end{enumerate}
困难在于，从形式上没有任何理由保证紧支撑推前 $\pi_!$ 保持反常性或不可约性。本章的关键是：将构造改写为 Fourier–Deligne 变换\cite{KL85}后，这个问题就能够处理。
\section{向量丛的 Fourier 变换}\label{sec:5.2}
设 $V\to S$ 为秩 $r$ 向量丛，$V^\vee\to S$ 为其对偶，$\langle-,-\rangle:V\times_S V^\vee\to\A^1$ 为配对。Fourier–Deligne 变换\cite{KL85}为
\[\mathsf F_{V/S}(\cK):=p_!^\vee(p^*\cK\otimes\langle-,-\rangle^*\cL_\psi)[r](r/2),\qquad\cK\in D_c^b(V).\]
采用这一归一化，它对于反常 $t$-结构是 $t$-正合的，并保持单反常层的不可约性。因此，问题是将定义 $\Aut'_E$ 的层识别为一个洁净延拓的 Fourier 变换。
\begin{transnote}{含糊点拨}
归一化 $[r](r/2)$ 的作用是把对偶的两支拉平：$p^\vee_!$ 与 $p^\vee_*$ 经 Verdier 对偶只差一个移位和扭转，取 $[r](r/2)$ 后 $\mathsf F_{V/S}$ 对反常 $t$-结构 $t$-正合。$r/2$ 可以是半整数，故系数域须含 $\sqrt q$，$\Qell$ 满足这一点。「保持单反常层的不可约性」不是算出来的，而是 $t$-正合加光滑性的形式后果：$t$-正合函子把单对象送到单对象（可能带平移）。[依据：原书 p.33–34 本段；并见 [KL85] 第 2 章的 $\mathsf F$ 归一化与 t-正合性。]
\end{transnote}
\section{\texorpdfstring{$G=\GL_2$}{G=GL₂} 的情形}\label{sec:5.3}
对 $G=\GL_2$，记
\[\overline{\Bun}'_2=\{0\longrightarrow\Omega_X\longrightarrow M\longrightarrow L\longrightarrow0,\quad L\in\Coh_1\}\]
为凝聚扩张叠。用于自守层的叠 $\Bun'_2$，是其中中间项 $M$ 局部自由的开子叠。令 $\Coh_1$ 为秩一凝聚层的叠，并令
\[\Coh_1^+:=\{(L,\cO_X\xrightarrow{s}L)\},\qquad\Coh'_1\subset\Coh_1^+\]
为 $s$ 在一般点单射的开子叠。在 $\Coh'_1$ 上，$s$ 的余核是挠层；令 $q:\Coh'_1\to\Tor$ 为这个商映射。因此，下文所用的层是 $q^*\cL_E$，而不是 $\cL_E$ 到 $\Coh'_1$ 的字面限制。

存在笛卡尔方块
\origdiagram{p033-square.pdf}{0.5692}
其中，给定扩张 $0\to\Omega_X\to M\to L\to0$ 以及 $\cO_X\to L$ 后，丛 $M'$ 是拉回。在 $\Coh_1$ 中满足 $H^1(X,L)=\Ext^1(\cO_X,L)=0$ 的开子叠上，$\Coh_1^+\to\Coh_1$ 是纤维为 $\Hom(\cO_X,L)$ 的向量丛。纤维为 $\Ext^1(L,\Omega_X)$ 的扩张叠是对偶向量丛，因为 Serre 对偶给出
\[\Ext^1(L,\Omega_X)\simeq\Hom(\cO_X,L)^\vee.\]
中间项局部自由的开子叠，是 $\Bun'_2$ 的相应部分。因此，必须在 $\Coh_1$ 的这个开子叠上工作，并且必须使用 $\Coh_1^+$ 而不是 $\Coh'_1$，因为零截面是向量丛的一部分。

Whittaker 映射 $\square\to\Bun_N\to\Ga$ 恰好是配对
\[\Hom(\cO_X,L)\times\Ext^1(L,\Omega_X)\longrightarrow\Ext^1(\cO_X,\Omega_X)\simeq k.\]
\srcpage{34}
因此，在 $\Coh_1$ 的这个开子叠上，
\[\Aut'_E\simeq\mathsf F(j_!q^*\cL_E),\qquad j:\Coh'_1\hookrightarrow\Coh_1^+.\]
若 $j_!=j_{!*}=j_*$，则 $j_!q^*\cL_E$ 为单反常层，因而其 Fourier 变换是不可约反常层。这至少在 $\Bun'_2$ 上具有向量丛描述的开部分证明了所需陈述。
\section{Drinfeld 关于 \texorpdfstring{$\GL_2$}{GL₂} 的洁净延拓定理}\label{sec:5.4}\index{洁净延拓|hyperpage}
现在限制到线丛 $\Bun_1\subset\Coh_1$。由于线丛之间的映射要么为零，要么为单射，在 $\Bun_1$ 上有
\[\Bun'_1\xrightarrow{j}\Bun_1^+\xleftarrow{i}\Bun_1,\]
其中 $i$ 为零截面，$j$ 为它的开补集。
\begin{theorem}{Drinfeld，洁净延拓；见\cite{FGKV98,KL85}}{5.4.1}
假设 $E$ 几何不可约且秩为 $2$。若 $d>4(g-1)$，则层 $j_!q^*\cL_E^d$ 在 $\Bun_1^d\subset\Coh_1$ 上洁净。
\end{theorem}
\begin{proof}

在 $\Bun'_1$ 上令 $\cF=q^*\cL_E^d$。$j_!\cF$ 的洁净性，是指自然映射 $j_!\cF\to j_*\cF$ 为同构。由于锥支撑于零截面，这等价于
\[i^*j_*\cF=0.\]
由 Verdier 对偶，只需对 $\cF$ 及其对偶证明紧支撑的边界消失。对于向量丛 $\Bun_1^+\to\Bun_1$ 上 $\Gm$ 的伸缩作用，收缩原理给出紧支撑形式
\[i^!j_!\cF\simeq\pi_!\cF,\]
其中 $\pi:\Bun'_1\to\Bun_1$ 为投影。模去叠性的 $\Gm$ 因子后，$\pi$ 是 Abel–Jacobi 映射
\[\operatorname{AJ}^d:X^{(d)}\longrightarrow\Pic^d.\]
因此，边界项由 Deligne 定理消失：对几何不可约且秩 $n>1$ 的 $E$，或对秩为 $1$ 且非平凡的 $E$，
\[d>n(2g-2)\quad\Longrightarrow\quad\operatorname{AJ}_!^d(E^{(d)})=0.\]
对 $n=2$，这一条件恰好是 $d>4(g-1)$。
\begin{transnote}{含糊点拨}
这是 Deligne 关于 Abel–Jacobi 的消失定理在秩 $n$ 的推广，也正是 \thref{5.7.2} 在 $k=1$ 时的特例：$\Av_E^d$ 在 $\Bun_1$ 上就是 $\operatorname{AJ}^d_!$ 沿 $E^{(d)}$ 的积分。$E^{(d)}$ 是 $E^{\boxtimes d}$ 的 $S_d$-不变部分在对称积 $X^{(d)}$ 上的下降；界 $d>n(2g-2)$ 随秩线性放大。$n>1$ 时必须设 $E$ 几何不可约，是为了排除 $E$ 含来自真抛物子群或真子局部系统的分量——那时 $E^{(d)}$ 混入低秩成分，消失一般不成立；$n=1$ 的「非平凡」是同一条件的退化形式，此时即经典结果。书中只引用不证明。[依据：原书 p.34–35；并见 [KL85] 的 Abel–Jacobi 消失定理。]
\end{transnote}
\end{proof}
\section{\texorpdfstring{$\GL_n$}{GLₙ} 的朴素迭代}\label{sec:5.5}
对一般 $n$，希望迭代上述 Fourier 变换。令 $\Coh_k$ 为秩 $k$ 凝聚层的叠。定义
\[\Coh_k^+:=\{(M_k,s:\Omega_X^{k-1}\to M_k)\},\qquad\Coh'_k\subset\Coh_k^+\]
为 $s$ 在一般点单射的开轨迹。若 $s$ 在一般点单射，则 $\operatorname{coker}(s)$ 的秩为 $k-1$，所以 $\Coh'_k$ 映到 $\Coh_{k-1}$。递归从沿商映射 $q_1:\Coh'_1\to\Tor$ 的拉回开始：
\[F_1:=q_1^*\cL_E,\]
并朴素地置
\[F_{k+1}:=\mathsf F(j_{k,!}F_k)|_{\Coh'_{k+1}},\qquad j_k:\Coh'_k\hookrightarrow\Coh_k^+,\]
其中将 $\Coh'_{k+1}$ 看作 $\Coh_k$ 上 Serre 对偶扩张丛的开子叠。最后，$F_n$ 限制到 $\Bun'_n$ 后，就是 $\Aut'_E$ 的 Fourier 描述。

不能简单地将所有中间的 $\Coh_k$ 限制到 $\Bun_k$。若对 $k<n$ 有 $C_k\subset\Bun_k$，则该构造只能看到紧化 Whittaker 叠中间商为向量丛的部分。
\srcpage{35}
这样会漏掉某些 $M_k\to M$ 具有挠余核的缺陷层。因此，正确的开子叠 $C_k\subset\Coh_k$ 必须比 $\Bun_k$ 更大。
\section{开子叠 \texorpdfstring{$C_k$}{Cₖ} 与两个所需条件}\label{sec:5.6}\index{归一化次数|hyperpage}
这一构造把 $C_k$ 的最终定义推迟了。这里使用的两个条件如下。设 $M_k\in C_k$，以 $M_k^{\mathrm{tf}}$ 表示它的无挠商。

首先，向量丛条件为
\[\Ext^1(\Omega_X^{k-1},M_k)\simeq\Ext^1(\Omega_X^{k-1},M_k^{\mathrm{tf}})=0.\]
它意味着 $\Coh_k^+\to\Coh_k$ 是纤维为 $\Hom(\Omega_X^{k-1},M_k)$ 的向量丛。同样的 Serre 对偶计算还给出 $\Hom(M_k,\Omega_X^k)=0$，所以对偶扩张叠没有自同构方向。它的线性对偶是扩张类构成的真正向量丛
\[\mathcal E_{k+1}\longrightarrow C_k,\qquad0\longrightarrow\Omega_X^k\longrightarrow M_{k+1}\longrightarrow M_k\longrightarrow0.\]
开子叠 $C'_{k+1}\subset\mathcal E_{k+1}$，是中间项属于所选开子叠 $C_{k+1}$ 且诱导截面 $\Omega_X^k\to M_{k+1}$ 在一般点单射的轨迹。其对偶性来自 Serre 对偶：
\[\Ext^1(M_k,\Omega_X^k)\simeq\Hom(\Omega_X^{k-1},M_k)^\vee.\]

其次，归一化次数必须足够大：
\[\ndeg(M_k^{\mathrm{tf}}):=\deg(M_k^{\mathrm{tf}})-\deg(\cO_X\oplus\Omega_X\oplus\cdots\oplus\Omega_X^{k-1})>nk(2g-2).\]
从这里起，即使省去上圈，所有次数也都指归一化次数 $\ndeg$。
\section{洁净性与平均消失}\label{sec:5.7}\index{平均函子|hyperpage}\index{消失定理|hyperpage}
在需要进行下一次 Fourier 变换的阶段，预期的洁净延拓陈述为：
\begin{theorem}{洁净性}{5.7.1}
对 $1\le k<n$，延拓 $j_{k,!}F_k$ 至少在 $\Bun_k\cap C_k\subset\Coh_k$ 上洁净。
\end{theorem}
若 $F_k$ 是 $C'_k$ 上的单反常层，则这一洁净延拓仍为单层。在 $C_k$ 的定义中加入额外条件后，同样的洁净性将在整个 $C_k$ 上成立。

洁净性的输入是 Gaitsgory 关于平均函子的消失定理\cite{Gai04,Gai16b}。令 $\operatorname{Mod}_k^d$ 为 Hecke 叠
\[\operatorname{Mod}_k^d=\{M\subset M',\ \length(M'/M)=d\},\]
具有映射
\origdiagram{p035-diagram.pdf}{0.3438}
其中 $\pi(M\subset M')=M'/M$。平均函子为
\[\Av_E^d(-):=\overrightarrow h_!(\overleftarrow h^*(-)\otimes\pi^*\cL_E^d)[dk](dk/2).\]
\srcpage{36}
\begin{theorem}{消失}{5.7.2}
若 $E$ 几何不可约且秩为 $n$，$k<n$，并且 $d>nk(2g-2)$，则
\[\Av_E^d=0.\]
更一般地，当 $\rank(E)>k$ 时，消失定理适用于 $\Bun_k$ 上的平均函子，其界为 $d>(2g-2)k\cdot\rank(E)$。
\begin{transnote}{含糊点拨}
两个界要合起来读：第一个界 $d>nk(2g-2)$ 里 $n=\rank E$、$k$ 是 $\Bun_k$ 的下标；第二个界 $d>(2g-2)k\rank E$ 把 $k=\rank E$ 的情形单列出来。因为 $\Av_E^d$ 的核是 $\pi^*\cL_E^d$，而 $\cL_E^d$ 的「体积」随 $\rank E$ 与 $d$ 同时增长，两条界只是同一增长率的不同写法。在整章证明里，这个定理只负责「挠部分为零」这一环，其余交给 Fourier 变换的 $t$-正合性；证明本身不在书中展开。[依据：原书 p.36 定理 5.7.2 及其后的说明；并见 [Gai04]、[Gai16b]。]
\end{transnote}
\end{theorem}

当 $k=1$ 时，这就是 Deligne 关于 Abel–Jacobi 的消失定理。\chref{6}用\thref{5.7.2}证明洁净性定理。
